\documentclass[11pt]{article}
\usepackage[T1]{fontenc}
\usepackage[utf8]{inputenc}
\usepackage{lmodern}
\usepackage{amsmath,amssymb,amsthm,mathtools}
\usepackage{graphicx}
\usepackage{microtype}
\usepackage[margin=1in]{geometry}
\usepackage[colorlinks=true,linkcolor=blue,citecolor=blue,urlcolor=blue]{hyperref}

\newtheorem{theorem}{Theorem}
\newtheorem{lemma}[theorem]{Lemma}
\newcommand{\Per}{\operatorname{Per}}
\newcommand{\Leb}{\mathcal L^1}
\newcommand{\lip}{\operatorname{lip}_a}

\title{An atomic-boundary minimizer with no closed representative\\
\large A counterexample to Question 5.2 of arXiv:2305.02891}
\author{Autonomous proof attempt, lane 6}
\date{13 August 2026}

\begin{document}
\maketitle

\begin{center}
\fbox{\parbox{0.90\textwidth}{\textbf{Status: candidate full counterexample,
likely valid.}  Question 5.2 has a negative answer for both functionals in
the general class of metric measure spaces stated by the source.}}
\end{center}

\section{The source question}

For a bounded domain $\Omega$ in a metric measure space
$(X,d,\mathfrak m)$, Koivu, Lu\v ci\'c, and Rajala consider
\[
 M_\lambda(A)=\Per(A)+\lambda\mathfrak m(\Omega\setminus A)
 \quad(A\subset\Omega)
\]
and
\[
 \widetilde M_\lambda(A)
 =\Per(A)+\lambda\mathfrak m(\Omega\mathbin\triangle A)
 \quad(A\subset X).
\]
Their Question 5.2 asks whether a Borel minimizer of either functional must
have a closed representative \cite[Question 5.2]{KLR}.  Here
\emph{representative} means equality up to an $\mathfrak m$-null set.  The
exact question appears on PDF page 17 of the source and is reproduced in
Figure~\ref{fig:question}.

\section{Proof intuition}

Put positive atoms exactly at the two boundary points of an interval.
Because convergence in $L^1(\mathfrak m)$ detects the atoms, every Lipschitz
approximation to the indicator of the open interval must be near zero at both
endpoints and near one somewhere inside.  It therefore pays at least one unit
of variation going up and one going down.  Thus the open interval has
perimeter two and no non-null competitor can have smaller perimeter.

For a fidelity weight $\lambda\ge2$, deleting the interval costs at least
$\lambda$, while including either endpoint in the symmetric-difference
problem also costs at least $\lambda$.  The open interval is consequently a
minimizer of both functionals.  But a closed set equal to the interval almost
everywhere would have to exclude the two positive-mass endpoints; closedness
then forces it to omit an entire one-sided neighborhood, of positive Lebesgue
measure.  Hence no closed representative exists.

\section{The counterexample}

Let
\[
 X=\mathbb R,\qquad d(s,t)=|s-t|,\qquad
 \mathfrak m=\Leb+\delta_0+\delta_1,
 \qquad \Omega=(0,1).
\]
This is a complete separable metric measure space, its Borel measure is finite
on bounded sets and has full support, and $\Omega$ is a bounded domain.  We
write $\Per(A)=|D\chi_A|(X)$ for the relaxed metric perimeter used in the
source.

\begin{lemma}\label{lem:perimeter}
\begin{enumerate}
\item $\Per(\Omega)=2$.
\item If $A\subset X$ is Borel, $0,1\notin A$, and
      $\Leb(A\cap(0,1))>0$, then $\Per(A)\ge2$.
\end{enumerate}
\end{lemma}

\begin{proof}
For the upper bound in (1), let $0<\varepsilon<1/4$ and choose the piecewise
linear function $u_\varepsilon$ which is zero outside
$(\varepsilon,1-\varepsilon)$, rises linearly from zero to one on
$[\varepsilon,2\varepsilon]$, equals one on
$[2\varepsilon,1-2\varepsilon]$, and falls linearly to zero on
$[1-2\varepsilon,1-\varepsilon]$.  Then
$u_\varepsilon\to\chi_\Omega$ in $L^1(\mathfrak m)$.  It is constant near
both atoms, so their contributions to
$\int_X\lip(u_\varepsilon)\,d\mathfrak m$ vanish, while the Lebesgue
contribution is exactly two.  Hence $\Per(\Omega)\le2$.

For the lower bound, consider any Lipschitz relaxation sequence $u_j$ for
$\chi_A$.  Composing with the truncation onto $[0,1]$ does not increase the
asymptotic slope, so assume $0\le u_j\le1$.  Since the measure has atoms at
$0$ and $1$ and these points are not in $A$,
\[
 u_j(0)\longrightarrow0,\qquad u_j(1)\longrightarrow0.
 \tag{1}
\]
If $\Leb(A\cap(0,1))>0$, convergence in the Lebesgue part of
$L^1(\mathfrak m)$ gives points $t_j\in A\cap(0,1)$ for which
$u_j(t_j)\to1$.  The fundamental theorem for Lipschitz functions and the
inequality $|u_j'|\le\lip(u_j)$ almost everywhere give
\begin{align*}
 \int_X\lip(u_j)\,d\mathfrak m
 &\ge \int_0^1\lip(u_j)\,dx\\
 &\ge |u_j(t_j)-u_j(0)|+|u_j(1)-u_j(t_j)|.
\end{align*}
Taking the lower limit and using (1) proves $\Per(A)\ge2$.  Applying this to
$A=\Omega$ gives the reverse inequality in (1).
\end{proof}

\begin{theorem}\label{thm:counterexample}
For every $\lambda\ge2$, the Borel set $E=\Omega$ is a minimizer of both
$M_\lambda$ among Borel subsets of $\Omega$ and
$\widetilde M_\lambda$ among Borel subsets of $X$.  Nevertheless, $E$ has no
closed representative.  Therefore Question 5.2 of \cite{KLR} has a negative
answer for both functionals.
\end{theorem}

\begin{proof}
First consider $M_\lambda$.  If $A\subset\Omega$ and $\Leb(A)>0$, then
Lemma~\ref{lem:perimeter} gives
\[
 M_\lambda(A)\ge\Per(A)\ge2.
\]
If $\Leb(A)=0$, then $A$ is $\mathfrak m$-null, because neither endpoint
belongs to $\Omega$.  Consequently
\[
 M_\lambda(A)=\lambda\mathfrak m(\Omega)=\lambda\ge2.
\]
On the other hand,
$M_\lambda(\Omega)=\Per(\Omega)=2$, so $\Omega$ is a minimizer.

Now consider an arbitrary Borel $A\subset X$ for
$\widetilde M_\lambda$.  If $0\in A$ or $1\in A$, then the corresponding
unit atom lies in $\Omega\triangle A$, and hence
\[
 \widetilde M_\lambda(A)\ge\lambda\ge2.
\]
If neither endpoint lies in $A$ and $\Leb(A\cap\Omega)>0$, then
Lemma~\ref{lem:perimeter} again yields
$\widetilde M_\lambda(A)\ge2$.  In the remaining case
$0,1\notin A$ and $\Leb(A\cap\Omega)=0$, the whole Lebesgue mass of $\Omega$
is missing, so
\[
 \widetilde M_\lambda(A)\ge
 \lambda\mathfrak m(\Omega\setminus A)=\lambda\ge2.
\]
Since $\widetilde M_\lambda(\Omega)=2$, the same $E=\Omega$ minimizes the
second functional.

Suppose finally that a closed $F\subset X$ were a representative of
$\Omega$.  Equality modulo $\mathfrak m$ forces $0,1\notin F$, since both
points have positive atomic mass and lie outside $\Omega$.  Since $F$ is
closed and $0\notin F$, some interval $(-r,r)$ is disjoint from $F$.
Then $(0,r)\subset\Omega\setminus F$ has positive Lebesgue measure,
contradicting $\mathfrak m(F\triangle\Omega)=0$.  Thus no closed
representative exists.
\end{proof}

\section{Verification and scope}

The proof uses only the source definitions.  The lower perimeter bound is
insensitive to behavior outside $[0,1]$: endpoint atoms force the two endpoint
values in every relaxation sequence, and one interior point with value
approaching one forces total variation at least two.  The explicit tent
sequence proves sharpness and has zero asymptotic slope at the atoms.

The example is deliberately outside the PI setting.  This is necessary for
the $\widetilde M_\lambda$ part, because the source already proves that
$\widetilde M_\lambda$ minimizers have closed representatives in PI-spaces.
Question 5.2, however, is stated for arbitrary metric measure spaces, exactly
the class containing the example.

A bounded novelty search covered the run indexes and the exact wording of
Question 5.2.  It also covered arXiv:2305.02891, its published IMRN version,
and the later arXiv:2503.15716 paper by overlapping authors.  No source
resolving Question 5.2 was found.  The IMRN version still labels it open.
This is a bounded search, not a definitive novelty certification.

\paragraph{Human-review recommendation.}
Review as a likely valid full counterexample.  The central check is the
relaxed-perimeter lower bound in Lemma~\ref{lem:perimeter}; all minimization
and representative claims then follow by a three-case comparison.

\clearpage
\begin{figure}[p]
  \centering
  \includegraphics[width=0.96\textwidth]{figures/open_problem_crop.png}
  \caption{Question 5.2 on PDF page 17 of arXiv:2305.02891.}
  \label{fig:question}
\end{figure}
\clearpage

\begin{thebibliography}{9}
\bibitem{KLR}
J. Koivu, D. Lu\v ci\'c, and T. Rajala,
\emph{Approximation by BV-extension sets via perimeter minimization in metric
spaces}, arXiv:2305.02891 (2023); Int. Math. Res. Not. IMRN 2024, no. 11,
9359--9375.
\end{thebibliography}

\end{document}
